\documentclass[10pt]{article}
\usepackage[margin=1in]{geometry}
\usepackage{amsmath}
\usepackage{graphicx}
\usepackage{authblk}

\title{Bridging measurements and simulations: A data assimilation approach to real time analysis}
\author{\underline{Ángel González Villatoro}, Mariachiara Gallia, David Rival}
\affil{Institut für Strömungsmechanik, Technische Universität Braunschweig, Germany }

\begin{document}
\maketitle

\begin{abstract}
Accurate reconstruction of wall-bounded flows from sparse measurements depends not 
only on the assimilation method and measured variables, but also critically on the 
spatial distribution of sensors. While previous data assimilation (DA) studies have 
successfully combined numerical simulations with experimental observations to improve 
reconstruction accuracy (Hayase \& Hayashi, 1997; Neeteson \& Rival, 2020), the question 
of where sensors should be placed remains central. Optimal sensor positioning has been 
shown to significantly enhance estimation quality while reducing instrumentation 
requirements (Mons \& Marquet, 2021).

In this work, we extend our previous analysis of variable selection by investigating 
the influence of sensor placement along the surface of a flat plate undergoing a sudden 
acceleration boundary-layer transition. High-resolution simulations provide the reference 
dataset, from which synthetic wall observations (temperature, pressure, shear stress) 
are sampled according to different spatial distributions. Strategies tested include 
uniform spacing, localized clustering near flow features, and adaptive placement based 
on wall-gradient indicators. These observations are assimilated into coarse-resolution 
simulations using lightweight DA schemes such as nudging and Kalman-based filters.

Performance is assessed in terms of reconstruction accuracy (RMSE), assimilation cost, 
and robustness to sparse configurations. Results show that sensors placed near the 
leading edge and in regions of strong shear provide the highest gains in reconstruction 
quality. Uniform distributions yield balanced improvements, while adaptive gradient-based 
placement achieves the best trade-off between accuracy and sensor economy.

These findings highlight the importance of optimized sensor placement for DA in 
wall-bounded flows, providing guidelines for integrating surface-based measurements into 
computational fluid dynamics workflows with minimal instrumentation requirements.
\end{abstract}

\end{document}

